%% CAPITULO 3 -- PROCEDIMENTOS METODOLOGICOS (esqueleto da dissertação)

\section{Procedimentos metodológicos}

A pesquisa segue o paradigma da \textit{Design Science Research} e se organiza
como um experimento computacional pareado. Sobre um corpus de documentos
institucionais públicos já reunido pelo autor, serão comparadas quatro
estratégias de contextualização (\textit{chunking} fixo com sobreposição,
recursivo, semântico e \textit{parent-child} sensível à estrutura), além de uma
ablação de metadados. Cada estratégia responderá ao mesmo conjunto de 60 a 100
perguntas de referência, incluindo perguntas sem resposta suficiente no corpus,
com índice vetorial exato e modelo gerador mantidos constantes, conforme o fluxo
da Figura~\ref{fig:fluxo-avaliacao}.

\begin{figure}[htb]
  \caption{Fluxo de avaliação das estratégias de contextualização documental}
  \label{fig:fluxo-avaliacao}
  \centering
  \fbox{\parbox{2.7cm}{\centering\small Documentos institucionais públicos}}
  $\rightarrow$
  \fbox{\parbox{2.7cm}{\centering\small Estratégia de contextualização}}
  $\rightarrow$
  \fbox{\parbox{2.7cm}{\centering\small Recuperação e geração da resposta}}
  $\rightarrow$
  \fbox{\parbox{2.7cm}{\centering\small Avaliação em 5 dimensões e IRD}}
  \fonte{Elaborada pelo autor (2026).}
\end{figure}

Por ser um instrumento novo, o IRD será validado antes da execução principal.
Os seis critérios passam primeiro por validação de conteúdo com um especialista
em Ciência da Informação; um piloto com pelo menos 30 respostas servirá apenas
para revisar o guia de codificação; em seguida, um estudo com cerca de 100
respostas codificadas por dois avaliadores independentes estimará o alfa de
Krippendorff por critério, e o menor valor será tomado como a confiabilidade do
índice. A validade convergente será verificada contra uma nota humana de
rastreabilidade atribuída por dois especialistas de outros órgãos públicos.
